%
% File: teach/elmagii/gallery/gallery-swedish.tex [plain TeX code]
% Github: https://github.com/elmagii/gallery/
% Last change: January 22, 2026
%
% Gallery of primary physicists who formed the theory of electomagnetism
% through history. These physicists and mathematicians essentially formed
% the theory which we are describing in the course ``Elektromagnetism II,
% 1TE626'', held 2022--2026 at Uppsala University, Sweden.
%
% Copyright (C) 2026, Fredrik Jonsson, under Gnu General Public
% License (GPL) v3. See the enclosed LICENSE for details.
%
% This program is free software: you can redistribute it and/or modify
% it under the terms of the GNU General Public License as published by
% the Free Software Foundation, either version 3 of the License, or
% (at your option) any later version.
%
% This program is distributed in the hope that it will be useful,
% but WITHOUT ANY WARRANTY; without even the implied warranty of
% MERCHANTABILITY or FITNESS FOR A PARTICULAR PURPOSE.  See the
% GNU General Public License for more details.
%
% You should have received a copy of the GNU General Public License
% along with this program.  If not, see <https://www.gnu.org/licenses/>.
%
\input macros/epsf.tex
\input macros/eplain.tex

\def\coursename{Elektromagnetism II}
\def\coursecode{1TE626}
\def\courseyear{2026}
\def\courserepo{https://github.com/hp35/elmagii/}
%-------------------- BEGIN OF LOCAL MACROS --------------------
\font\ninerm=cmr9
\font\tenssbx=cmssbx10
\font\twelvesc=cmcsc10 at 12 truept
\def\initgallery{
  \noindent~\vskip-30pt\hskip-44pt%
    {\epsfxsize=90pt\epsfbox{../lect-01/macros/UU_logo_color.eps}}
  \vskip-64pt\hfill\vbox{
      \hbox{{\it \coursename, \coursecode\ (\courseyear)}}
      \hbox{{\it Document Revision \today}}
      \hbox{{\it \courserepo}}}\vskip 36pt\noindent}
%--------------------- END OF LOCAL MACROS ---------------------

%
% Define language constants. We are here only concerned with Swedish or English.
%
\newcount\swedish \swedish=0
\newcount\english \english=1

%
% Set the current language to be used throughout.
%
\newcount\lang
\lang=\swedish  % Change to \english for English

%
% Define the 'boxit' macro from D.E. Knuths "The TeXbook, Exercise 21.3.
% Set the parameter \boxedtrue or \boxedfalse in order to have boxed
% portrait generation.
%
\def\boxit#1{\vbox{\hrule\hbox{\vrule\kern1pt
  \vbox{\kern1pt#1\kern1pt}\kern1pt\vrule}\hrule}}
\newif\ifboxed
\boxedtrue   % \boxedtrue or \boxedfalse

\def\presection#1{%
  \noindent{\font\twelvecmbx=cmbx12\twelvecmbx{#1}}%
  \smallskip\noindent}%

\initgallery
\input src/portrait.tex
\input src/gallery-src.tex

\bye
